Définir la stratégie SSI & P09, P02, P17 & Niveau 1 étayé par la note d'orientation 2024 approuvée par la DG, des priorités annuelles et un référent (P09). Niveau 2 indéterminé car aucune mesure qualitative n'est documentée et le compte rendu semestriel reste oral, sans trace écrite fournie. Niveau 3 non atteint, faute de processus commun documenté (P09). Revue annuelle de la PSSI sans PV depuis deux ans (P17). & Tableau de bord qualitatif semestriel et comptes rendus écrits validés par la DG, intégrés au lot A01 (pilote RSSI). Preuve : deux PV de revue datés. \\
Gérer les risques SSI & P10, P06 & Risques recensés ponctuellement lors de projets récents, avec des mesures décidées (P10). Niveau 2 non atteint : ni calendrier de revue, ni méthode réutilisable, ni registre consolidé approuvé. L'analyse du prestataire de 2025 reste isolée. & Méthode simple, registre consolidé et revue trimestrielle (A01). Preuve : registre approuvé par la DG et PV de revue. \\
Gérer les règles SSI & P11, P17, T06 & Niveau 2 étayé par une PSSI approuvée, des règles écrites, une application planifiée, des contrôles qualitatifs consignés et l'information des managers (P11). Niveau 3 non atteint, les textes ne forment pas un processus commun. Fragilités : revue annuelle sans PV (P17) et aucune cible de délai de correctifs (T06). & PSSI révisée avec délais de correctifs par criticité, diffusée à tout le périmètre (A03). Preuve : PSSI v2 approuvée, PV de revue, taux de sensibilisation. \\
Superviser la SSI & P12, P06, T08 & Outil de journalisation et alertes actifs, interventions à la demande (P12), d'où le niveau 1. Niveau 2 non atteint : aucune campagne de contrôle planifiée, aucun reporting régulier ni procédure commune (P12). AL-1048 fermée automatiquement sans ticket (T08). & Plan de contrôle annuel, procédure de traitement des alertes et reporting mensuel (A02). Preuve : échantillon de tickets d'alerte qualifiés et rapports mensuels. \\
Concevoir les mesures SSI & P13, T02 & Besoins identifiés sur certains projets, fiche type, responsable technique désigné, revues planifiées pour les projets majeurs avec compte rendu écrit remonté au responsable (P13). Le compte rendu écrit des revues est retenu comme mesure qualitative. La compétence est retenue, comme pour les autres processus, sur la base d'un rôle désigné chargé de la sécurité, à confirmer en entretien. Niveau 3 non atteint, la démarche n'est pas globale. La règle provisoire FW-17 (T02) l'illustre sans le prouver. & Fiche de besoins SSI obligatoire pour tout projet ou changement significatif (préparée en régie, A04 différée). Preuve : fiches sur un échantillon de projets. \\
Réaliser les mesures SSI & P14, T02 & Niveau 1 étayé par des déploiements préparés par des techniciens compétents (P14). Comptes rendus d'installation et de validation, suivi régulier et information de la DSI couvrent quatre exigences du niveau 2. Niveau 2 indéterminé car la planification n'est pas démontrée, P14 indiquant seulement que les déploiements sont « souvent préparés ». Règles FW-17 et FW-31 sans justification ni validation (T02). & Plan de déploiement trimestriel des mesures et modèle commun de recette, en régie (pilote DSI). Preuve : plans datés et PV de recette sur trois déploiements. \\
Exploiter les mesures SSI & P15, P17, T03, T04, T06, T09 & Sauvegardes planifiées, comptes rendus hebdomadaires, astreinte, procédure de restauration connue et bilan qualitatif mensuel à la DSI (P15). Niveau 3 non atteint. Fragilités : 2 PV de test sur 9 services (T09), MFA sur 6 comptes privilégiés sur 10 (T04), correctifs en retard sans exception (T06), déclarations de la DG non démontrées et de la DSI contredites (P17). & Standardiser restauration, correctifs et revue des comptes (A05, A08, A09). Preuve : PV de restauration avec RTO et RPO mesurés, export MFA, SLA de correctifs. \\
